% !TeX root = ../sections/02_exercises.tex
\subsection{1-B-2.}
\begin{exercise}
	以下の等式を $\epsilon N$ 論法で証明せよ。
	\[
		\lim_{n\to\infty}
		\sqrt{n}\left(\sqrt{n+1}-\sqrt{n}\right)
		=
		\frac{1}{2},
	\]
	\[
		\lim_{n\to\infty}
		\frac{n^2+2n}{n+1}
		=
		\infty.
	\]
\end{exercise}
\begin{background}
	この問題では，有限の極限と正の無限大への発散を，
	それぞれ定義に従って示す。

	\begin{definition}[数列の収束]
		数列 $\{a_n\}$ が実数 $\alpha$ に収束するとは，
		任意の $\epsilon>0$ に対して，ある自然数 $N$ が存在して，
		$n\geq N$ ならば
		\[
			|a_n-\alpha|<\epsilon
		\]
		となることをいう。
	\end{definition}

	\begin{definition}[正の無限大への発散]
		数列 $\{a_n\}$ が正の無限大に発散するとは，
		任意の $M>0$ に対して，ある自然数 $N$ が存在して，
		$n\geq N$ ならば
		\[
			a_n>M
		\]
		となることをいう。
	\end{definition}

	\begin{remark}
		第一の極限では，有理化を用いて
		\[
			\sqrt{n}(\sqrt{n+1}-\sqrt{n})
			=
			\frac{\sqrt{n}}{\sqrt{n+1}+\sqrt{n}}
			=
			\frac{1}{\sqrt{1+\frac{1}{n}}+1}
		\]
		と変形する。

		第二の極限では，多項式の割り算により
		\[
			\frac{n^2+2n}{n+1}
			=
			n+1-\frac{1}{n+1}
		\]
		と変形できる。
	\end{remark}
\end{background}
\begin{strategy}
	第一の極限では，まず有理化を行う。

	\[
		\sqrt{n}(\sqrt{n+1}-\sqrt{n})
		=
		\frac{1}{\sqrt{1+\frac{1}{n}}+1}
	\]

	したがって，示すべきことは
	\[
		\frac{1}{\sqrt{1+\frac{1}{n}}+1}
		\to
		\frac{1}{2}
	\]
	である。

	差を評価すると，
	\[
		\left|
			\frac{1}{\sqrt{1+\frac{1}{n}}+1}
			-
			\frac{1}{2}
		\right|
	\]
	を $1/n$ の形で上から押さえることができる。
	したがって，任意の $\epsilon>0$ に対し，
	$n$ を十分大きく取ればこの差を $\epsilon$ より小さくできる。

	第二の極限では，
	\[
		\frac{n^2+2n}{n+1}
		=
		n+1-\frac{1}{n+1}
	\]
	と変形する。

	この式は $n$ が大きくなるといくらでも大きくなる。
	実際，
	\[
		n+1-\frac{1}{n+1}\geq n
	\]
	であるから，任意の $M>0$ に対して $N>M$ となる自然数 $N$ を取れば，
	$n\geq N$ のとき
	\[
		\frac{n^2+2n}{n+1}>M
	\]
	が従う。
\end{strategy}